% Source: https://artofproblemsolving.com/community/c3h1941022_area_formula_for_a_parallelogram_using_diagonals
% Stable identity: aops-topic-1941022-q01
Let the parallelogram be $ABCD$ with diagonals $AC=d_1$ and $BD=d_2$ meeting at $O$. The diagonals of a parallelogram bisect each other, so $OA=OC=\dfrac{d_1}{2}$ and $OB=OD=\dfrac{d_2}{2}$.

\begin{asy}
usepackage("amsmath");
size(6.5cm);
defaultpen(fontsize(10pt));
pair A=(0,0),B=(6,0),D=(2,3),C=B+D;pair O=(A+C)/2;
draw(A--B--C--D--cycle,linewidth(1));
draw(A--C,linewidth(0.7));draw(B--D,linewidth(0.7));
draw(arc(O,0.6,degrees(C-O),degrees(D-O)),linewidth(0.6));
label("$\theta$",O+0.95*dir((degrees(C-O)+degrees(D-O))/2));
label("$A$",A,SW);label("$B$",B,SE);label("$C$",C,NE);label("$D$",D,NW);label("$O$",O,S+0.5*E);
label("$\dfrac{d_1}{2}$",(6.15,1.55));label("$\dfrac{d_2}{2}$",(O+D)/2,W+0.4*S);
label("\mbox{}",(min(currentpicture,user=true)+max(currentpicture,user=true))/2,black); // keeps the text after the figure black when LaTeX typesets the labels inline
\end{asy}

The diagonals cut the parallelogram into triangles $OAB$, $OBC$, $OCD$, $ODA$. Each has two sides $\dfrac{d_1}{2}$ and $\dfrac{d_2}{2}$ with included angle $\theta$ or $180^\circ-\theta$, and $\sin(180^\circ-\theta)=\sin\theta$. So each has area $\dfrac{1}{2}\cdot\dfrac{d_1}{2}\cdot\dfrac{d_2}{2}\sin\theta=\dfrac{d_1d_2\sin\theta}{8}$, and the total area is
\[
4\cdot\dfrac{d_1d_2\sin\theta}{8}=\dfrac{1}{2}d_1d_2\sin\theta .
\]
For $d_1=10$, $d_2=8$, $\theta=30^\circ$: the area is $\dfrac{1}{2}\cdot10\cdot8\cdot\dfrac{1}{2}=\boxed{20}$.
